\sectionguard
\section*{Appendix: Scope and Qualifications}

\begin{studyframe}
\textbf{Edition and use.} This companion studies the recurring
\latin{Dominica Duodecima post Pentecosten}, II class, in the 1962 Vatican
typical Roman Missal. It is not an official liturgical book, homily, civil-date
assembly sheet, or ruling on present authorization. Occurrence, precedence,
commemorations, and particular calendars remain outside the reusable
formulary.

\textbf{Text and translation.} All ten Latin elements, their boundaries,
references, rank, Credo, and Trinity Preface were visually collated against the
identified facsimile. Scripture English is the registered public-domain
Douay--Rheims/Challoner; oration English is the anonymous Cummiskey 1861
witness. Adapted chants are accompanied by witness-boundary notices rather
than project translations. The 1962 scan itself is not redistributed; the
focused Latin texts are pre-existing liturgical material whose rights analysis
is recorded in the shared inventory.

\textbf{Research boundary.} Complete biblical contexts, principal Greek and
Latin patristic expositions, representative medieval Doctors, later saints,
Catholic doctrinal controls, sacramentary transmission, and targeted cultural
afterlives were checked within the matrix in \texttt{research/scope.md}. The
survey is broad and claim-driven, not exhaustive. Textual observation,
documented reception, source-grounded synthesis, and exploratory proposal are
kept distinct. The enumerated historical witnesses do not establish a single
compiler's cross-proper intent or the universal novelty of an editorial
proposal.

\textbf{Review.} The Latin, historical English, claim-level sources,
bibliography, bindings, and component classifications were audited. Both
editions were built through settled passes, their logs checked, and the local
web source conversion validated. These are internal source and production
facts, not ecclesiastical approval or an installed-publication claim. Full
records: \nolinkurl{propers/verified.md} and
\nolinkurl{research/scope.md}.
\end{studyframe}
